% =============================================================================
% Project Warrant -- Appendix draft: Limitations and Dual-Use Considerations
%
% Source: adapted from planning/SPRINT-PLAN.md \S9 (gitignored, local-only;
% quoted in full in issue #74), updated with five caveats accumulated after
% \S9's original draft was written: the ledger freeze at v1.2/125 rows, the
% Gate G1 blind-recode reliability check, and the lw-01/lw-02 stress-test
% retry. This file is meant to replace the placeholder body of
% \S\ref{app:limitations} in report/report.tex (lines 141-155); it is kept as
% a standalone draft here per issue #74's requested output path and can be
% pasted in wholesale once the surrounding report is assembled.
% =============================================================================

\section{Limitations and Dual-Use Considerations}
\label{app:limitations}

\subsection*{Limitations}

This audit's conclusions are bounded by six limitations, each of which would
change how the headline metrics should be read if ignored.

\paragraph{Single-coder ledger, blinded re-code only.} The 125-row claim
ledger was built and coded by a single researcher, not by independent
multi-rater coding across the whole ledger. The reliability check that exists
is a 25-row blind re-code (Rater~A's original P1.1 verdicts held out and
compared against a cold, independently re-coded pass over the same 25 rows),
scored with Cohen's~$\kappa$ per degradation condition. That check found
$\kappa = 0.7788$ (D1, opacity) and $\kappa = 0.7794$ (D3, adversarial
forgery), both above the 0.6 bar, but $\kappa = 0.5690$ on D2
(unfaithfulness) -- below it. The D2 codebook rule was subsequently fixed and
the whole ledger's D2 column was recoded against the corrected rule before
freeze. This resolves the specific rule ambiguity the recode surfaced; it
does not substitute for full independent multi-rater coding across all 125
rows, and the remaining 100 rows outside the 25-row subsample were never
independently re-coded by a second rater under either the original or the
fixed rule.

\paragraph{Corroboration taken at face value.} Claims are coded as
\emph{reported}. A claim can be scored \texttt{SURVIVES} under D1 (opacity)
because a source \emph{asserts} the existence of corroborating telemetry,
logs, or a second channel that this project has no independent means to
inspect. This project cannot verify that an asserted corroborating channel
actually exists, was actually consulted, or would actually hold up if
produced. This is the ceiling on M2 (record survival under opacity, 90.4\%
overall) in particular: a claim scored as surviving opacity is a claim whose
\emph{stated} evidentiary chain does not require chain-of-thought, not a
claim this project has confirmed would survive an actual forensic
examination. We restate this here as a named limitation on the audit's
method, not only as a caveat embedded in the Results narrative.

\paragraph{Channel-biased claim population.} The claim population is drawn
from what the public record chose to publish -- incident reports, blog posts,
and third-party investigations that already made editorial choices about
which evidence to disclose, quote, or summarize. Channels that were never
written up publicly at all cannot appear in this ledger, and channels that
labs chose to describe only at a category level (rather than quoting
directly) are underrepresented in the ledger's degree of technical detail
regardless of how load-bearing they actually were. The ledger therefore
measures channel dependence \emph{within the published record}, not within
the full evidentiary record the investigating institutions actually held.

\paragraph{Degradation verdicts are analytic, not empirical.} The D1/D2/D3
verdicts (SURVIVES / DEGRADED / COLLAPSES) are judgments about whether a
claim's stated evidentiary warrant would hold up under a stress condition --
opacity, unfaithfulness, or adversarial forgery -- reasoned from the claim
and its cited channel(s) as reported. They are not empirical measurements of
what an investigator with real access to raw logs, model weights, or
infrastructure could actually recover, reconstruct, or independently verify
under any of these conditions. No degradation condition in this project was
tested against real systems; all three are applied as analytic overlays on
the claim ledger.

\paragraph{Single-incident calibration.} The codebook, the three degradation
conditions, and the five headline metrics were all built against one
incident. Nothing in this project establishes that the same channel codes,
verifiability tiers, or degradation rules would classify claims from a
different incident the same way, or that the specific fragile-channel
finding (M5: 78.4\% of all 125 claims rest on a codebook-fragile primary
channel) generalizes beyond this one case. The instrument's transfer to
other incidents is untested.

\paragraph{Small-N metrics.} Two of the headline results rest on thin
denominators and should be read as directional rather than precise. M1 (CoT
monopoly, 11.1\%) is computed over only 9 T3 (intent) claims in the entire
ledger; a single row moving to or from the sole-C1 condition would shift the
reported figure by roughly 11 percentage points. M2's per-claim-type
breakdown for T4 (counterfactual claims, 100\% survival under D1) rests on
only 2 claims total -- two data points cannot confirm or refute the sprint
plan's own prediction that T3 and T4 claims would both fail to survive
opacity; the honest reading is that T4 did not collapse in this ledger's two
counterfactual claims, not that the plan's prediction about T4 was
confirmed. Both caveats are documented in \texttt{docs/headline-metrics.md}
and are restated here as standing limitations on how much confidence M1 and
the T4 slice of M2 can bear.

\paragraph{Residual lw-01/lw-02 category-level gap.} The ledger was frozen
once at v1.1 (116 rows) and reopened a single, disclosed time to add 9 rows
(\texttt{C352}--\texttt{C360}) recovered from sources \texttt{lw-01} and
\texttt{lw-02}, bringing it to v1.2 (125 rows). That recovery pulled M3 and
M5 toward less dramatic values than the pre-recovery figures (M3: 9.5\%
$\rightarrow$ 8.8\%; M5: 84.5\%/50.9\% $\rightarrow$ 78.4\%/47.2\%), which is
itself evidence the original numbers were not inflated by the omission in
the other direction. This is a separate and narrower point from that
directional finding, however: the 9 recovered rows were extracted and coded
at a category-only level of description -- section headings and topic
identification only, never the underlying technical content, commands, or
payloads themselves. The fuller technical detail present in
\texttt{lw-01}/\texttt{lw-02}'s own source material is still not reflected
anywhere in the ledger. What this project reports about those two sources'
contribution to the record is therefore a coarser summary than what the
sources themselves actually contain, and that gap should not be read as
closed by the 9-row addition -- only partially narrowed.

\subsection*{Dual-Use Considerations}

This audit names which evidentiary channels are load-bearing for the
public's understanding of this incident. That is also, unavoidably, a map of
which channels an adversary -- or an institution with an incentive to
control the narrative of a future incident -- would most want to suppress,
withhold, or forge next time. We take this risk seriously and adopt two
mitigations, both already scoped in the sprint plan before this project's
data collection began.

First, the analysis is kept at the level of \emph{channel classes} (raw
chain-of-thought, inter-agent message content, tool-call/action traces, lab
self-assertion, and so on) rather than at the level of any specific,
reproducible tamper or spoofing technique. Per rule R-06, no note field in
the ledger, no discussion in this report, and no supporting document in this
repository reproduces a forgery or tampering mechanism at a level of detail
that would let a reader carry it out. Where the public source material
documents a successful tool-call-spoofing technique used during the
incident, this project logs its existence and category only -- consistent
with how every D3 (adversarial forgery) note in the ledger is written, and
with how the excerpt inventories treat comparable technique-level content
throughout. We state explicitly, per the sprint plan's own instruction on
this point: this project declines to reproduce the documented spoofing
technique, even though it is described in a public source available to us.
Publishing a channel-class ranking is a different act from publishing a
recipe, and this project confines itself to the former.

Second, the constructive half of this project -- the v0.1 minimum forensic
standard (MRFM) -- is written so that its defensive value exceeds its
offensive value. Its clauses tell a defending institution which channels
need redundancy, retention, and attestation; they do not tell an attacker
anything about specific channels' weaknesses that could not already be
inferred from the same published incident reports this audit draws on. A
channel ranking derived entirely from public retrospectives is not new
information handed to a capable adversary -- the labs' own reports already
disclose which channels existed and how they behaved under stress. What this
project adds is a systematic account of which channels the \emph{public
record's own claims} depend on, framed toward the question of what an
institution should preserve, corroborate, and attest to before the next
incident, not toward what a future intruder should target.
